\introduction % Структурный элемент: ВВЕДЕНИЕ

Распределённые системы являются основой современных вычислительных платформ,
облачных сервисов, систем хранения данных и инфраструктуры обработки задач.
Именно в таких системах особенно остро проявляются проблемы согласованности,
отказоустойчивости и корректного взаимодействия независимых узлов. Даже
небольшая ошибка в алгоритме координации может приводить к расхождению
состояний реплик, потере результатов вычислений, некорректному выбору лидера
или нарушению гарантий безопасности. Сложность анализа подобных ошибок
обусловлена тем, что поведение распределённой системы формируется в условиях
асинхронности, недетерминированной доставки сообщений, конкурентного
выполнения и возможных сбоев отдельных процессов и сетевых соединений.

Для повышения надёжности распределённых алгоритмов всё шире применяются
формальные методы. Одним из наиболее известных и практически значимых средств
формального описания является язык спецификаций TLA+,
предложенный Л. Лэмпортом \cite{lamport2002specifying}. Он позволяет задавать
систему как множество состояний и допустимых переходов между ними, а также
формулировать инварианты и свойства корректности, которые должны сохраняться
при любом выполнении. Средство проверки моделей TLC, входящее в экосистему
TLA+, делает возможным автоматическое исследование пространства состояний
спецификации и выявление нарушений сформулированных свойств.

Однако наличие корректной спецификации ещё не означает корректность реальной
программной реализации. Между формальной моделью и исходным кодом существует
естественный разрыв: спецификация описывает поведение на абстрактном уровне,
тогда как реализация включает сетевое взаимодействие, многопоточность,
внутренние буферы, особенности библиотек и инфраструктурные детали. Именно на
этом этапе могут появляться расхождения, которые не видны при проверке одной
лишь модели. Поэтому важной практической задачей становится не только
верификация спецификации, но и проверка того, что исполняемая программа
действительно следует этой спецификации.

В данной работе рассматривается подход к валидации реализации распределённого
алгоритма консенсуса Raft по формальной спецификации TLA+ на основе анализа
трасс исполнения. Для этого реализация инструментируется таким образом, чтобы
во время работы фиксировать события, соответствующие изменениям переменных
спецификации. Полученные трассы затем сопоставляются со спецификацией при
помощи TLC через специальную trace-обёртку. Такой подход позволяет проверять не
абстрактную модель отдельно от программы, а реальное поведение запущенной
системы, воспроизведённое в виде последовательности наблюдаемых переходов.

\subsection*{Актуальность темы исследования}

Актуальность темы определяется тем, что традиционные методы тестирования
распределённых систем ограниченно подходят для обнаружения ошибок, связанных с
редкими межпоточными сценариями, особенностями сетевого взаимодействия и
переходами между состояниями протокола. Модульные и интеграционные тесты обычно
проверяют отдельные заранее предусмотренные случаи, тогда как множество
фактически возможных сценариев в распределённой системе значительно шире.
Особенно трудно воспроизводимыми являются ошибки, возникающие при специфическом
порядке прихода сообщений, повторной доставке, смене лидера или конкуренции
между параллельными действиями узлов.

Формальная спецификация в TLA+ позволяет строго описать корректное поведение
алгоритма, однако без сопоставления со следами реального исполнения остаётся
открытым вопрос, действительно ли программа реализует именно эту модель.
Следовательно, необходим метод, который связывает спецификацию и код на уровне
наблюдаемого поведения системы. Перспективным решением является модельная
валидация трасс исполнения, при которой последовательность событий,
зафиксированных в программе, проверяется на соответствие множеству переходов,
разрешённых формальной спецификацией \cite{howard2011tracechecking}. Такой
подход позволяет обнаруживать не только дефекты реализации, но и ошибки
инструментирования, расхождения в уровне атомарности действий и слишком сильные
или, наоборот, неполные предположения в самой спецификации.

Практическая значимость такого подхода особенно высока для алгоритмов
консенсуса, поскольку именно они лежат в основе отказоустойчивых сервисов,
реплицированных хранилищ и систем координации. Алгоритм Raft широко применяется
как более понятная альтернатива Paxos и используется в ряде промышленных
систем. Поэтому разработка и апробация метода проверки соответствия реализации
Raft его формальной спецификации представляет как теоретический, так и
прикладной интерес. Подобные подходы к использованию формальных методов для
анализа реальных распределённых систем уже доказали свою эффективность в
индустрии \cite{newcombe2015aws}.

\subsection*{Цель и задачи работы}

Целью данной работы является разработка и экспериментальная апробация метода
проверки соответствия реализации распределённого алгоритма консенсуса Raft его
формальной спецификации на языке TLA+ на основе трасс исполнения программы.

Для достижения поставленной цели были сформулированы следующие задачи:

\begin{enumerate}
    \item Разработать инструменты подготовки, объединения и проверки трасс
исполнения относительно спецификации с использованием TLC.
    \item Провести экспериментальную проверку реализации системы распределенных
вычислений, сопоставить реальные трассы со спецификацией и проанализировать
обнаруженные расхождения.
\end{enumerate}
